\documentclass[11pt]{article}
\usepackage[a4paper,margin=25mm]{geometry}
\usepackage{amsmath,amssymb,amsthm,mathtools,booktabs,array,xcolor}
\usepackage[T1]{fontenc}\usepackage{lmodern}
\usepackage[colorlinks=true,linkcolor=teal,citecolor=teal,urlcolor=teal]{hyperref}
\hypersetup{pdftitle={Bicomplex Signal Realization and Fixed-Boundary Point Extensions},pdfauthor={Akari Hayami (Jian-Yu Huang)},pdfsubject={Local research continuation 0.02}}
\newtheorem{theorem}{Theorem}[section]
\newtheorem{proposition}[theorem]{Proposition}
\newtheorem{lemma}[theorem]{Lemma}
\newtheorem{corollary}[theorem]{Corollary}
\theoremstyle{remark}\newtheorem{remark}[theorem]{Remark}
\newcommand{\R}{\mathbb R}\newcommand{\C}{\mathbb C}
\newcommand{\Dom}{\operatorname{Dom}}\newcommand{\dd}{\,d}
\newcommand{\norm}[1]{\lVert#1\rVert}
\DeclareMathOperator{\rank}{rank}
\title{Bicomplex Signal Realization\\and Fixed-Boundary Point Extensions}
\author{Akari Hayami (Jian-Yu Huang)}
\date{1 October 2026\\\small Research continuation 0.02; local working draft}
\begin{document}\maketitle
\begin{abstract}
We give an explicit unit-power bicomplex state family whose phase-sensitive
quadratic readout reproduces the ruled signal map. The required phase coupling
is stated as model data. We also prove that three fixed Hermitian quadratic
observables on normalized two-component pure states cannot reproduce this
whole map. For the actual induced Whitney metric, a singular coordinate change
reduces the scalar Dirichlet form to a uniformly elliptic divergence-form
problem with bounded density. Interior H\"older estimates then give continuous,
graph-bounded values at both interior cross-caps. Consequently the corrected
point restriction of the fixed Dirichlet operator has deficiency indices
$(2,2)$, an actual resolvent boundary triple and a $U(2)$ extension family.
Its Green vectors are defined by the actual resolvent; explicit logarithmic
normalizations and singular Dirac transmission remain separate questions.
\end{abstract}

\section{An explicit signal realization}\label{sec:realization}
We use the same completed source labels and map as the verified v13 draft:
\[
 S(t,u)=\bigl(c+2u(1-c),(1-u)s,uQ\bigr),\quad
 c=\cos t,\ s=\sin t,\ Q=\sqrt{c(2-c)},
\]
where $-\pi/2\leq t\leq\pi/2$ and $0\leq u\leq1$.
The bicomplex idempotent coordinates are $z=z_+e_++z_-e_-$;
the chosen Hilbert power is $|z_+|^2+|z_-|^2$.
Define the two unit phasors
\[
 \zeta_+(t)=e^{it/2},\qquad
 \zeta_-(t)=\sqrt{1-c/2}+i\sqrt{c/2}.
\]
Thus $\zeta_+^2=c+is$ and $\zeta_-^2=1-c+iQ$.
These phases encode the selected orthogonal-circle coupling. Power
conservation alone does not select that coupling.

\begin{theorem}[Quadrature realization]\label{thm:realization}
The continuous state family and fixed real quadratic readout
\begin{align*}
 \Phi(t,u)&=\sqrt{1-u}\,\zeta_+(t)e_+
                 +\sqrt u\,\zeta_-(t)e_-,\\
 P(z)&=\bigl(\Re(z_+^2)+|z_-|^2+\Re(z_-^2),
                   \Im(z_+^2),\Im(z_-^2)\bigr)
\end{align*}
satisfy $|\Phi_+|^2=1-u$, $|\Phi_-|^2=u$ and $P\circ\Phi=S$
on the entire physical source. The lift is smooth where $0<u<1$ in
the completed source charts.
\end{theorem}
\begin{proof}
Both phasors have modulus one. Their squares give
$P\Phi=((1-u)c+u(2-c),(1-u)s,uQ)$, which is $S$.
The readout measures the real and imaginary quadratures of the component
squares, together with one DC power term. It is specified on the ambient
algebra before a source point is substituted.
At the angular endpoints use $r=Q$, $C=1-\sqrt{1-r^2}$ and
$t=\varepsilon\arccos C$. The imaginary component of $\zeta_-$ is
$r/\sqrt{2(2-C)}$ for $r\geq0$ and has a smooth signed extension;
its real component is $\sqrt{1-C/2}$. The first phasor is smooth in $C$.
This verifies smoothness in the actual endpoint charts. Square-root power
weights are smooth in the open ruling interval and continuous at its ends.
\end{proof}

The construction retains a common phase reference: multiplying both
components by $e^{i\alpha}$ rotates their squares by $e^{2i\alpha}$.
It is a classical phase-sensitive signal readout, rather than a claim about
three quantum expectation values.

\begin{proposition}[Power-coordinate boundary obstruction]\label{prop:powerboundary}
No $C^1$ lift in the original coordinate $u$ can satisfy $|z_-(t,u)|^2=u$
up to $u=0$. Similarly $|z_+(t,u)|^2=1-u$ prevents a $C^1$ lift at $u=1$.
Writing $u=\sin^2v$ gives a smooth lift on $0\leq v\leq\pi/2$,
but this substitution is not a boundary diffeomorphism of the original source.
\end{proposition}
\begin{proof}
At $u=0$ the component is zero. If it were $C^1$, the one-sided derivative
of its squared modulus would be
$2\Re(\overline{z_-}\,\partial_u z_-)=0$, contradicting the derivative of $u$.
The other endpoint is identical. In the power-angle coordinate the weights
are $\cos v$ and $\sin v$; the derivative of $\sin^2v$ vanishes at both ends.
\end{proof}

\begin{proposition}[Three Hermitian observables do not realize this map]\label{prop:hermitian}
There are no three fixed Hermitian $2\times2$ matrices $H_j$ and normalized
pure-state family $z(t,u)\in\C^2$ such that
$S_j(t,u)=z(t,u)^*H_jz(t,u)$ throughout the physical source.
\end{proposition}
\begin{proof}
The Bloch vector of a normalized two-component pure state lies on the unit
sphere. Expanding $H_j$ in the identity and Pauli matrices gives
$P_H(z)=a+Mr(z)$ with a real $3\times3$ matrix $M$;
see also the pure-state distinction in \cite{bloch}.
If $\rank M\leq2$, its image is planar. But the four actual values
$(1,0,0),(0,1,0),(0,-1,0),(1,0,1)$ of $S$ have affine determinant $2$.
If $M$ is invertible, its image of the sphere is an ellipsoid surface,
which contains no nonconstant straight segment. The actual ruling
$S(0,u)=(1,0,u)$ is such a segment. Both cases are impossible.
\end{proof}
The proposition concerns normalized pure states and fixed Hermitian
observables. It does not exclude mixed states, additional components, or
the phase-sensitive readout already constructed.

\section{Resolving the intrinsic scalar metric}\label{sec:resolved}
For the standard cross-cap $F(x,y)=(x,xy,y^2)$ set $R=\sqrt{x^2+y^2}$.
Its actual pullback metric and a new coordinate map are
\[
 g_0=dx^2+(y\,dx+x\,dy)^2+4y^2dy^2,
 \qquad \Psi(x,y)=(X,V)=(x,yR).
\]
The coordinate change is singular at zero. It preserves the completed
source point and its neighbourhood topology.

\begin{lemma}[Uniform metric comparison]\label{lem:metric}
The map $\Psi$ is a homeomorphism of $\R^2$ and a smooth diffeomorphism
off zero. For $0<R\leq1/2$,
\[
 \frac1{16}\Psi^*(dX^2+dV^2)\ \leq\ g_0\ \leq\
 8\Psi^*(dX^2+dV^2).
\]
Also $X^2+V^2=|F(x,y)|^2$.
\end{lemma}
\begin{proof}
For fixed $x$, the function $y\mapsto y\sqrt{x^2+y^2}$ is strictly
increasing onto $\R$. Its inverse is continuous and obeys
\[
 y^2=\frac{\sqrt{X^4+4V^2}-X^2}{2},\qquad \operatorname{sgn}y=\operatorname{sgn}V.
\]
Off zero the Jacobian is $(x^2+2y^2)/R>0$. Continuity at zero follows
from the inverse formula. The radius identity follows by expansion.
For a tangent vector $(\xi,\eta)$ put $h=\xi^2+R^2\eta^2$ and
$E=g_0((\xi,\eta),(\xi,\eta))$. The exact positive decompositions
\begin{align*}
 E-\tfrac12h&=\tfrac12(2y\xi+x\eta)^2
       +(\tfrac12-y^2)\xi^2+\tfrac72y^2\eta^2,\\
 4h-E&=(y\xi-x\eta)^2+(3-2y^2)\xi^2+2x^2\eta^2
\end{align*}
give $h/2\leq E\leq4h$.
Now $dV=a\,dx+b\,dy$, where $a=xy/R$, $b=(x^2+2y^2)/R$.
Here $|a|\leq R/2\leq1/4$ and $1\leq b/R\leq2$.
Consequently $\xi^2+(a\xi+b\eta)^2\leq8h$.
Conversely $R\eta=(R/b)((a\xi+b\eta)-a\xi)$ gives
$h\leq2(\xi^2+(a\xi+b\eta)^2)$. Combining these estimates proves
the displayed comparison.
\end{proof}

Grieser already proved that Whitney-umbrella metrics become quasi-Euclidean
under a singular coordinate change \cite[Section 1]{grieser}. We give the
explicit comparison here to track its scalar PDE hypotheses; the coordinate
method is not claimed as a new discovery.

\begin{corollary}[Application to the actual ruled germs]\label{cor:actualmetric}
Each of the actual interior germs $P_\pm$ admits a source homeomorphism,
smooth off its singular label, in which its metric tensor $G$ and density
$m=\sqrt{\det G}$ are bounded above and below by positive constants.
The divergence matrix $a=mG^{-1}$ is bounded, measurable and uniformly elliptic.
\end{corollary}
\begin{proof}
The verified adapted determinants in v13 show that each actual germ is a
smooth cross-cap. Thus an actual source diffeomorphism $h$ and an ambient
target diffeomorphism $L$ give $S=L\circ F\circ h$, after translations.
Shrink the neighbourhood so the singular values of $DL$ lie between positive
constants $\ell$ and $M$. Then $\ell^2g_0\leq F^*(L^*g_{\R^3})\leq M^2g_0$.
Apply Lemma~\ref{lem:metric} to $\Psi\circ h$. Bounds on the eigenvalues of
$G$ imply the asserted density and divergence-matrix bounds. This is
metric comparison through an actual target map, not an assertion that $L$
is an isometry or that it preserves sheet parity.
\end{proof}

\section{Actual graph-bounded point traces}\label{sec:trace}
Let $\mathcal H=L^2(D_{\mathrm{reg}},\dd A_g)$ and let $q_D$ be the
Friedrichs form of the actual scalar metric, with the regular outer Dirichlet
boundary fixed. Its self-adjoint operator is $H_F$.
The all-five-point capacity proof, compact form embedding and positive
Dirichlet gap are the proved v13 inputs. Only the two interior labels are
used for the point restriction below.

\begin{theorem}[Actual point evaluation]\label{thm:trace}
Every $f\in\Dom H_F$ has a unique continuous representative near $P_\pm$
in the completed source. In the resolved coordinates it is locally H\"older
continuous, and
\[
 \tau f=(f(P_+),f(P_-)),\qquad
 \norm{\tau f}\leq C(\norm f_{\mathcal H}+\norm{H_Ff}_{\mathcal H}).
\]
The trace $\tau:\Dom H_F\to\C^2$ is surjective and its kernel is dense in
$\mathcal H$.
\end{theorem}
\begin{proof}
On a resolved coordinate disk the form norm is equivalent to $H^1$ and
the Hilbert norm to ordinary $L^2$, by Corollary~\ref{cor:actualmetric}.
A point has zero $H^1$ capacity in dimension two: cutoffs on a logarithmic
annulus have energy tending to zero. Thus the local completed form space
is $H^1$ on the whole disk. This also follows by localizing the original
Friedrichs core and completing it in the equivalent norm.
Writing $b=H_Ff$, its weak equation becomes
\[
 -\operatorname{div}(a\nabla f)=mb.
\]
It holds across zero. Indeed tests on the whole disk can be approximated
in $H^1$ by tests avoiding zero, and both the form pairing and the $L^2$
right-hand side are continuous in that norm. There is no unaccounted delta
term in this equation for a Friedrichs-domain function.

The coefficients are real symmetric, measurable and uniformly elliptic;
$m$ is bounded; and $mb\in L^2$, with $2>\dim(D)/2=1$.
The interior inhomogeneous De Giorgi--Nash estimate gives, on a smaller disk,
\[
 \norm f_{C^\alpha}\leq C\bigl(\norm f_{L^2}+\norm{mb}_{L^2}\bigr).
\]
This is the standard external scalar theorem, for example
\cite[Theorem 8.24]{gt}; the homogeneous estimate and $L^p$, $p>n/2$,
inhomogeneous version are also stated in \cite[Questions 52,58]{silvestre}.
For a local formulation one can subtract the zero-boundary solution of the
inhomogeneous equation on the disk and apply the homogeneous estimate to
the remainder. Apply the theorem to real and imaginary parts separately.
This yields a continuous value and the claimed graph estimate.
Two continuous representatives equal almost everywhere must agree, since
the resolved density is bounded below. The coordinate homeomorphism
then supplies a unique value at the actual source label.

Choose two disjoint source cutoffs equal to one near the respective point
and zero near all outer boundaries and the other singularities. Capacity
zero puts them in the form domain; their Laplacians are smooth and supported
on regular annuli, so they are in $\Dom H_F$. Their traces are the two
standard basis vectors, proving surjectivity. Finally the original punctured
interior smooth compact core lies in $\ker\tau$ and is dense in $\mathcal H$.
\end{proof}

\begin{corollary}[The corrected two-point operator]\label{cor:indices}
The actual operator
\[
 A=H_F|_{\ker\tau}
\]
is closed, densely defined, symmetric, strictly positive and has deficiency
indices $(2,2)$. Its Friedrichs extension is $H_F$ and all its self-adjoint
extensions are parametrized by $U(2)$.
\end{corollary}
\begin{proof}
Graph continuity makes $\ker\tau$ graph closed. The remaining basic properties
are inherited from $H_F$ and Theorem~\ref{thm:trace}. For nonreal $z$,
$\tau(H_F-z)^{-1}$ is a bounded surjection to $\C^2$ with kernel
$\operatorname{ran}(A-z)$. Hence this closed range has codimension two,
so $\dim\ker(A^*-\bar z)=2$. Von Neumann's unitary classification applies.
The punctured form core is contained in $\Dom A\subset\Dom H_F$;
its form closure is exactly $q_D$. Therefore the Friedrichs extension is $H_F$.
\end{proof}
This construction is the fixed-boundary point restriction used in regular
point-interaction theory, e.g.\ \cite{point}, with its graph-trace hypothesis
now verified for the actual Whitney metric. The old graph closure of the
interior compact core remains a different operator with infinitely many
outer-boundary deficiency modes; its historical $(2,2)$ statement stays false.

\section{Actual Green vectors and boundary pairing}\label{sec:green}
The point traces allow an exhaustive domain description without guessing a
logarithmic coefficient. Fix a real $\eta<0$ and set
\[
 R_z=(H_F-z)^{-1},\qquad G_z=[\tau R_{\bar z}]^*:\C^2\to\mathcal H.
\]
All adjoints use a conjugate-linear first inner product.

\begin{theorem}[Resolvent boundary triple]\label{thm:triple}
The actual adjoint domain decomposes as
\[
 \Dom A^*=\Dom H_F\dotplus G_\eta\C^2,
 \qquad A^*(u+G_\eta a)=H_Fu+\eta G_\eta a.
\]
With $\Gamma_0(u+G_\eta a)=a$ and $\Gamma_1(u+G_\eta a)=\tau u$,
\begin{align*}
 \langle A^*f,g\rangle-\langle f,A^*g\rangle
 =\langle\Gamma_1f,\Gamma_0g\rangle
       -\langle\Gamma_0f,\Gamma_1g\rangle.
\end{align*}
The pair $(\Gamma_0,\Gamma_1)$ is surjective onto $\C^2\oplus\C^2$.
Every self-adjoint extension is specified by
\[
 (U-I)\Gamma_1f+i(U+I)\Gamma_0f=0,\qquad U\in U(2).
\]
In this convention $U=I$ gives $H_F$.
\end{theorem}
\begin{proof}
The bounded map $T_\eta=\tau R_\eta$ is onto; its adjoint is injective and
has range $\operatorname{ran}(A-\eta)^\perp=\ker(A^*-\eta)$.
For $f\in\Dom A^*$ put $u=R_\eta(A^*-\eta)f$. Then $f-u$ belongs to
that kernel. The intersection with $\Dom H_F$ is zero since $\eta$ is in
the resolvent set of $H_F$. This proves the unique decomposition and action.
For $f=u+G_\eta a$, $g=v+G_\eta b$, the $u,v$ and $a,b$ terms cancel.
The cross terms use
\[
 \langle(H_F-\eta)u,G_\eta b\rangle=\langle\tau u,b\rangle
\]
and its conjugate, giving the stated Green form. Given $(a,b)$, take $u$
with $\tau u=b$ and set $f=u+G_\eta a$. The kernel of both traces is
$\Dom A$. Self-adjoint restrictions correspond to maximal isotropic planes
of this finite nondegenerate Green form. The Cayley transform gives exactly
the displayed unitary equation, including planes that are not graphs of
a Hermitian matrix.
\end{proof}

\begin{corollary}[Actual finite-rank resolvent formula]\label{cor:resolvent}
Let $B=B^*$ and impose $\Gamma_1=B\Gamma_0$. With
\[
 M(z)=(z-\eta)\tau R_zG_\eta,
\]
the resolvent, wherever both sides are defined, is
\[
 (H_B-z)^{-1}=R_z+G_z(B-M(z))^{-1}G_{\bar z}^*.
\]
Thus the difference from the fixed Dirichlet resolvent has rank at most two.
\end{corollary}
\begin{proof}
The resolvent identity gives $G_z=G_\eta+(z-\eta)R_zG_\eta$;
in particular $\Gamma_0G_za=a$ and $\Gamma_1G_za=M(z)a$.
A solution of $(A^*-z)f=h$ has the unique form $f=R_zh+G_za$.
Its boundary condition becomes $(B-M(z))a=\tau R_zh=G_{\bar z}^*h$.
Solving this finite system gives the formula.
\end{proof}

\begin{corollary}[Compact extensions and the actual Krein kernel]\label{cor:krein}
Every self-adjoint extension of $A$ has compact resolvent. Its
Krein--von Neumann extension $A_K$ has a two-dimensional kernel and satisfies
\[
 \Dom A_K=\Dom A\dotplus G_0\C^2,
 \qquad \Gamma_1f=M(0)\Gamma_0f.
\]
\end{corollary}
\begin{proof}
For nonreal $z$, the difference of the resolvents of any extension and $H_F$
has range in $\ker(A^*-z)=G_z\C^2$, since both solve the same adjoint
equation. Hence this difference has rank at most two. The already proved
compactness of $R_z$ gives compactness of every extension resolvent.
Strict positivity and the positive-operator Krein theorem
\cite[Theorem 2.1, equation (2.8)]{krein} give
$\Dom A_K=\Dom A\dotplus\ker A^*$.
Because $0$ belongs to the resolvent set of $H_F$, the same trace-resolvent
argument at zero yields $\ker A^*=G_0\C^2$, of dimension two.
Finally $\Gamma_1G_0=M(0)\Gamma_0G_0$ and both traces vanish on $\Dom A$,
which identifies its boundary plane.
\end{proof}
This is the Krein extension of the corrected point restriction $A$.
The old compact-core operator still has an infinite-dimensional Krein kernel.

The vectors $G_ze_\pm$ are actual graph-dual point-source Green vectors.
The theorem excludes additional defect directions for this specified operator
$A$; it does not identify their coefficients with those of a particular
$\log\rho$ ansatz. Changing the reference $\eta$ changes the regular
boundary coordinate and Weyl normalization, not the actual extension family.
An explicit local logarithmic expansion remains useful for geometric
interpretation and quantitative couplings.

\section{What this establishes for the Dirac problem}\label{sec:dirac}
The scalar theorem does not require derivatives of the resolved coefficient
matrix. A first-order spin operator does, through its connection, so its
domain cannot be inferred from scalar uniform ellipticity.

\begin{proposition}[Actual curvature integrability]\label{prop:curvature}
At an actual smooth cross-cap germ, the Gaussian curvature belongs locally
to $L^p(\dd A_g)$ for $1\leq p<3/2$. On the standard germ,
\[
 K=-\frac{4y^2}{(x^2+4y^2+4y^4)^2}.
\]
In a source disk of radius $\varepsilon$, the absolute curvature integral
is at most $C\varepsilon$.
\end{proposition}
\begin{proof}
For the actual map in smooth adapted source coordinates, all second
derivatives are bounded. The coefficients of its second fundamental form
are consequently bounded, because the unit normal has length one.
Metric comparison with the standard germ gives $\det g\asymp R^2$ and
$\dd A_g\leq C R\,dx\,dy$. The Gauss formula therefore gives
$|K|\leq C/R^2$. Polar integration bounds the $p$th curvature integral by
$C\int_0^\varepsilon R^{2-2p}\dd R$, finite precisely in the stated sufficient
range. For $p=1$ this is $C\varepsilon$.
For the standard map direct first- and second-fundamental-form calculation
gives the displayed exact formula.
\end{proof}

This supplies an actual integrable spin-curvature input: the curvature of
the two-dimensional spin connection is proportional to $K\dd A_g$,
and tensoring by the flat observation sign line adds no local curvature.
It does not by itself provide a bounded conformal factor, a controlled
connection gauge, or the weighted graph-domain comparison needed at the
singular cut. Those estimates, the phase-line fibre identification and a
Clifford-compatible transmission map remain the next Dirac obligations.
The one-channel-per-chirality exact-cone calculation retains its model scope.

\section{Evidence and relation to the verified checkpoint}\label{sec:evidence}
\begin{center}\small
\begin{tabular}{p{.27\linewidth}p{.64\linewidth}}\toprule
Result & Evidence and scope\\\midrule
Signal realization & Explicit global continuous lift, interior/completed-chart smoothness and endpoint obstruction; quadratic identities replayed exactly.\\
Metric comparison & Written homeomorphism and positive estimates; standard metric energies checked independently and in Lean. Grieser's earlier coordinate result is credited.\\
Point traces & Actual metric comparison and form-domain argument, followed by the external scalar H\"older theorem; not Lean formalized.\\
Extension domains & Written Hilbert-space proof for $A=H_F|_{\ker\tau}$, actual Green vectors and finite-rank formula; no claimed logarithmic parametrix.\\
Dirac progress & Actual curvature bound and exact standard curvature formula; singular spin-domain and transmission classification remain open.\\\bottomrule
\end{tabular}
\end{center}
This continuation adds results to the locally verified v13/companion 0.01
checkpoint. Its historical PDF, sealed package, manifests and receipts are
unchanged. Exact algebra and the listed Lean theorems do not certify the
external PDE regularity theorem or the full operator-domain argument.
The explicit signal family is new supplied model data, not recovered
historical source. Novelty of the complete application requires further
literature comparison; no physical capacity, particle or chirality prediction
is established by these results.

\begin{thebibliography}{9}
\bibitem{grieser} D. Grieser, \emph{Quasiisometry of singular metrics},
Houston J. Math. \textbf{28} (2002), 741--752.
\bibitem{gt} D. Gilbarg and N. S. Trudinger,
\emph{Elliptic Partial Differential Equations of Second Order},
Springer, Classics in Mathematics (2001), Chapter 8, Theorem 8.24.
\href{https://doi.org/10.1007/978-3-642-61798-0}{Publisher edition}.
\bibitem{silvestre} L. Silvestre, \emph{Boot camp -- Problem set},
29 September 2017, Sections 4.1 and 6, Questions 52 and 58.
\url{https://math.uchicago.edu/~luis/preprints/bootcamp.pdf}.
\bibitem{point} D. Noja and F. Raso Stoia,
\emph{The Dirichlet Laplacian with a point interaction on unbounded
Lipschitz domains}, arXiv:2607.04349v1 (2026), Proposition 2.1 and Theorem 2.2.
\bibitem{krein} M. S. Ashbaugh, F. Gesztesy, M. Mitrea, R. Shterenberg
and G. Teschl, \emph{The Krein--von Neumann extension and its connection
to an abstract buckling problem}, arXiv:0907.1439v2 (2010), Section 2.
\bibitem{bloch} K. Szyma\'nski and K. \.Zyczkowski,
\emph{Geometric and algebraic origins of additive uncertainty relations},
arXiv:1804.06191v1 (2018), Section II.
\end{thebibliography}
\end{document}
